
%
% !BIB TS-program = biber
% !BIB program = biber
% !TeX spellcheck = en_US
% !TeX program = lualatex
%
% support document for version: v 3.00  Feb 2026   Volker RW Schaa
%   
%
\documentclass[%              % paper size definition not needed as JACoW page size is set
               %boxit,        % show page boundaries to check whether paper is inside correct margins
               %refpage       % references placed on separate page
               %biblatex,     % BibLaTeX is to be used
               %nospread,     % flushend option: do not fill with whitespace to balance columns
               %hyphens,      % allow \url to hyphenate at "-" (hyphens)
               %luatex,       % use LuaLaTeX to process the file
               ]{jacow}

\usepackage[english]{babel}

\listfiles

%%
%% Lengths for the spaces in the title
%%   \setlength\titleblockstartskip{..}  %before title, default 3pt
%%   \setlength\titleblockmiddleskip{..} %between title + author, default 1em
%%   \setlength\titleblockendskip{..}    %afterauthor, default 1em

\begin{document}

%%%%% Title - Author block generated from indico contribution details %%%%%
%%%%%%%%%%%%%%%%%%%% Do not modify  %%%%%%%%%%%%%%%%%%%%%%%%%
%%%%% If you need to modify the content of this block, do it by modifying the  %%%%%
\title{SPIN COHERENCE TIME OPTIMIZATION BY MAPPING METHOD}

\author{S.~Kolokolchikov\thanks{sergey.bell13@gmail.com}, A.~Aksentyev, A.~Melnikov, P.~Palamarchuka, Y.~Senichev, Russian Academy of Sciences, Moscow, Russia}

\maketitle

%%%%% End of Title - Author block %%%%%

%%%%% You can modify the abstract to add LaTeX commands. However, please make sure that it is similar to the one in indico %%%%%
\begin{abstract}
A sufficiently long Spin Coherence Time (SCT) is a fundamental requirement for precision measurements of the electric dipole moment (EDM) of charged particles in storage rings. SCT optimization is performed using three sextupole families that suppress second-order effects arising from emittance and momentum spread. These sextupoles are placed at the maxima of the beta functions in both planes and in regions with significant horizontal dispersion. The study examines the capability of the mapping method to simultaneously suppress spin decoherence and achieve zero chromaticity in both transverse planes in magnetic, electrostatic, and combined-function structures.
\end{abstract}

\section{MANUSCRIPTS}

This short document demonstrates the use of the JACoW formatting style in \LaTeX, highlighting the template's key features. It is intended to help authors quickly create a correctly formatted JACoW manuscript in \LaTeX. Authors are expected to be familiar with the contents of the "JACoW Style Guide" document (\verb|jacow_latex_style_guide|), which is available on the JACoW website.

\subsection{Headings}

Title and section headings are automatically converted to uppercase, therefore the author must ensure 
that all units and acronyms with mixed case writings (keV, MeV, GHz, µs, SwissFEL, SuperKEKB, 
FCC-hh) are correctly typeset (see use of macro \verb|\NoCaseChange{text}| in above title).
The plural of proper nouns and acronyms in uppercase has to be a LOWERCASE “s” (e.\,g. 
LINAC\textbf{s}, PLM\textbf{s}). 

\subsection{Figures and Tables}

Figures and tables should be placed as close as possible to where they are mentioned. 
Each figure and table must be numbered in ascending order (1, 2, 3, etc.) throughout the paper. 
The lettering in figures and tables must be large enough to be reproduced clearly. 

When referring to a figure from within the text, the
convention is to use the abbreviated form (e.\,g., Fig.~\ref{fig:example})
\emph{unless} the reference is at the start of the sentence. Reference to a
table, however, is never abbreviated (e.\,g., Table~\ref{tab:example}).
As a matter of style, authors are advised not to begin any sentence with an abbreviation or a number.

Table captions are typically composed as a heading, with the initial letters of principal words capitalized and without a period at the end. A description of the table contents should be confined to the text of the paper and not in the table caption. 
\begin{table}[!hbt]
   \centering
   \caption{Table Heading}
   \begin{tabular}{lcc}
       \toprule
       \textbf{Column 1} & \textbf{Column 2} & \textbf{Column 3} \\ 
                         & \textbf{[mm]}     & \textbf{[GeV]}	 \\
       \midrule
          Attribute A    & \num{100.2}       & \num{20.3}        \\
          Attribute B    & \num{137.0}       & \num{~9.1}        \\
       \bottomrule
   \end{tabular}
   \label{tab:example}
\end{table}

\subsection{Equations}

%If the command \qty is generating an error, uncomment the line below this one:
%\ifdefined\qty  \else \newcommand{\qty}{\SI} \fi

Equations with numbers do not need to be referenced (the number can be suppressed using \verb|\nonumber|):
\begin{equation}\label{eq:label}
    C_B=\frac{q^3}{3\epsilon_{0}mc} = \qty{3.54}{\micro eV/T}\ .
\end{equation}

\subsection{Units}

Units should be written using \verb|\qty{value}{unit}| or \verb|\num{value}| of the “siunitx” package as shown in Eq.~\eqref{eq:label}. Some examples are: \qty{3}{keV}, \qty{100}{kW}, \qty{7}{\um}.

\subsection{References}

In the interest of promoting complete and uniform citations, 
the IEEE Editorial Style for Transactions and Journals has been adopted. 
References~\cite{ref:jacow-paper, ref:journal-paper} exemplify this style. 
When citing a periodical, use the ISO~4 standard for abbreviating journal names. 
Please refer to Annex~C of the “Style Guide for JACoW Conference Papers” for examples. 
Authors are responsible for paying attention to the details of the style to ensure their references are 
accurate, and properly formatted. 
Each paper published in recent JACoW proceedings has been assigned a digital object identifier (DOI). Including a DOI in a reference facilitates importing references into bibliographic databases and is strongly recommended. 
Furthermore, hyperlinks to DOIs are required, while links to non-DOI URLs are prohibited. 
JACoW conferences that stipulate a page limit typically exempt the reference section from this requirement.

\section{CONCLUSION}

This short document illustrates the use of the JACoW format style, which is described in detail in the Style Guide. 

%\begin{thebibliography}{99}   	% Use for  10-99  references
\begin{thebibliography}{9} 		% Use for 1-9 references
	
	\bibitem{ref:jacow-paper}
		A. Alpha and B. T. Beta, 
		“Novel techniques for future TeV electron accelerators”, 
		in \textit{Proc. IPAC’23}, Venice, Italy, May 2023, pp.~57-59. 
	
	\bibitem{ref:journal-paper}
		A. Alpha \textit{et al.}, 
		“A new 12\,GHz electron linear accelerator”, 
		
\end{thebibliography}

\end{document}